\section{The Second ``C'': Adaptive Unary-Path Compression}\label{sec:unary-path}


In this section, we address the space inefficiency incurred by unary paths
(especially suffix unary paths) and introduce the second "C" of \csquared: the
adaptive unary-path \underline{C}ompression scheme (denoted with
\compressionscheme). \rev{The \compressionscheme optimization enables all tries
  to access unary-path compression in the data. Specifically, it combines
  Marisa's recursion~\cite{marisa} with the Fast Static Symbol Table
  (\fsst)~\cite{fsst} as the tail container. Furthermore, it \emph{adaptively}
  chooses the number of levels of recursion for the Marisa trie based on
  per-dataset space savings.

  We evaluate \csquared tries against baselines at the same recursion level;
  since CoCo-trie and FST originally lack recursion, \compressionscheme
  changes only their tail container.}

\para{Locality issues in the data} \rev{Long unary paths increase topology size
  and waste space through redundant label sequences.} \newrev{Without proper
  compression, suffix paths can account for about 80\% of succinct-trie space,
  as shown in~\tabref{unary-wiki-log}.}

\newrev{Prior work uses two orthogonal compression methods (detailed
  in~\secref{preliminaries}): PDT~\cite{pdt} uses \repair~\cite{larsson2000off},
  while Marisa~\cite{marisa} uses recursion, as detailed
  in~\secref{preliminaries}. Neither method dominates: on the
  \texttt{log}~\cite{log-dataset} dataset, recursion and \repair achieve
  compression ratios of 16.4$\times$ and 7.5$\times$, respectively; on
  \texttt{xml}~\cite{text-collection}, the ratios are 3.9$\times$ and
  5.0$\times$.} %\todo{update ratios}


\para{Choice of tail container} Although Marisa originally uses a sorted
  tail container, recursion is agnostic to the tail container; more advanced containers such as approximate
  \repair~\cite{pdt} and \fsst~\cite{fsst} can improve both space and query
  time.

  We limit our \csquared tries to the \fsst tail container because it captures almost all (within \fsstvsrepair) of the compression benefits of
  approximate \repair with $12.8\times$--$21.6\times$ faster build times. 
  %We omit the relevant
  %data due to space constraints but will include it in the full \newmymarginpar{R3 D7}\newrev{(i.e., ArXiv)} version. 
  As
  shown in~\tabref{c2-full} in~\secref{experiments}, choosing \fsst as the tail
  container always achieves a better compression ratio (up to \fsstmaxspacesavings) with minimal impact on the query time (.96-1.13$\times$)
  over the sorted container.
  %across all tested succinct tries.


  %Furthermore, changing the tail container to \fsst improves query time by
  %$1.12\times$ on average compared to sorted, because \bv operations in the
  %topology \new{and/or intra-node lookup} make up the majority of the time in
  %each query.

   % \new{As a result, \compressionscheme limits the choice of tail container to
  % \fsst to achieve a good balance between build time and compression ratio. We
  % enable the user to choose to optimize more for build time or compression
  % ratio by choosing the sorted or approximate \repair tail containers,
  % respectively. We note that the gains in either dimension from other tail
  % containers are small relative to \fsst.}

\begin{figure}[t]
    \begin{minipage}[h]{.23\textwidth}
      \includegraphics[width=\textwidth,page=9,trim=14cm 7.6cm 15cm
      7.5cm,clip]{visualization.pdf}
        \caption{Suffix path compression for \csquared-FST and \csquared-CoCo.}
        \label{fig:adaptive-compression-fst-example}
    \end{minipage}
    \hspace{0.1cm}
    \begin{minipage}[h]{.23\textwidth}
      \includegraphics[width=\textwidth,page=10,trim=14cm 7.8cm 14cm
      7.2cm,clip]{visualization.pdf}
        \caption{A \csquared-CoCo suffix path compression example. Stored keys:
          cash, camp, cell, crash.} % Exact lookup for "cell" would fail.}
        \label{fig:adaptive-compression-coco-example}
    \end{minipage}
\end{figure}


\para{Integration with succinct tries} \rev{First, we enable \bvdesign-FST,
  \bvdesign-CoCo, and \bvdesign-Marisa to access recursion and compressed tail
  containers. We focus the discussion on %integration with
  FST and CoCo, as these previously could not access unary-path compression.}

\figref{adaptive-compression-fst-example} illustrates how we containerize the
storage of suffix paths in FST. Specifically, each leaf node has an \islink bit indicating whether its
suffix path is moved to the next container. A suffix link at position $i$ can thus be uniquely identified by
\var{LinkId(i) = IsLink.rank$_1$(LeafID(i))}. When a lookup operation reaches a
leaf node, we follow the \var{IsLink} bitvector to check the remaining suffixes
in the container (if any).

The integration with \rev{\bvdesign-CoCo} is \rev{similar}\mymarginpar{R2 D15}
to \rev{\bvdesign-FST}, but \rev{CoCo requires special care for correctness: its
  encoding schemes support only exact
  lookup}.~\figref{adaptive-compression-coco-example} shows the issue. When
looking up the string ``cell,'' the CoCo-trie would search for ``cel'' in the
root \rev{because the root node has depth 3}. Since ``cel'' is not a key of the
root, the lookup would fail, even if ``cell'' exists in the trie.
% This error happens because all encoding schemes used by CoCo-trie only support
% exact lookup, while the target entry could be a prefix of current key.
% Given that keys are sorted within each macro node,
Instead, we search for the lower bound of the target key rather than an exact match. If the lower bound is a prefix of the target key, we trace the link to
check the remaining suffixes.

\rev{\para{Adaptive recursion depth}} We enable all of the \rev{succinct} tries
studied in this paper to access unary-path compression, but focus on the case of
the Marisa trie and show \new{how to achieve better space-time tradeoffs by
  combining recursion with \fsst.}

Using recursion in the other tries (FST and CoCo) trades query performance for
space savings. Therefore, we limit non-Marisa tries to no recursion in our evaluation,
but expose it as a user option for different space-time tradeoffs. We will include the
data on the effect of recursion in non-Marisa tries in the full version.

\new{The original Marisa-trie exposes the maximum number of recursions as a
  parameter to the user, but offers no per-dataset guidance on the space-query
  tradeoff. As the original Marisa trie documentation notes~\cite{marisa-code}
  and \secref{experiments} confirms, recursion has a serious negative impact on
  query performance, so we only want to continue recursion when the space
  reduction is significant.

  We \emph{adaptively select} the number of recursion levels that achieve
  the best time-space tradeoff per dataset. To balance
  compression ratio against build/query time, we stop recursion when the space
  savings would be less than some $\epsilon = 0.1$\mymarginpar{R1 D3, R2 D5}
  (relative to trie size). However, the user can set $\epsilon$ to suit
  their use case.}

\mymarginpar{R1 W3, R1 D3, R2 D1, R2 D4} To efficiently determine the space
usage at each level of recursion, we adopt \fsst's fast estimation scheme,
encoding a subset of the data to approximate the compression ratio.  In
practice, this estimate is within 10\% of the true compression ratio.


\para{Other locality optimizations} \new{For the Marisa trie, compressed tail
  containers degrade lookup performance in nodes with many links. To address
  this issue, we store the branching label (i.e., the first label of each unary
  path) in the label vector, enabling in-place intra-node search accelerated by
  SIMD~\cite{pdt}.}

\mymarginpar{R1 D3, R2 D1, R2 D3, R3 O4}\rev{To further improve locality, we
  store unary paths smaller than a link in-place, avoiding the index overhead of
  moving them to the next level. The number of bits needed to store a link
  depends on the size of the trie (i.e., if a trie has $n$ nodes, we need
  $\lg(n)$ bits to store each link). Short unary paths may be further
  compressed by moving them to the next level, but the space savings is limited
  by the short path length relative to the link cost.}


%%% Local Variables:
%%% mode: latex
%%% TeX-master: "main"
%%% End:
